\chapter{Lista 5.}
\begin{problem}[1]

	Załóżmy, że grupa H działa przez automorfizmy na grupie N. Dowieść, że elementem odwrotnym do $\left(n,h\right)$ w otrzymanym iloczynie półprostym $N\rtimes H$, jest element $\left( h^{-1}\cdot n^{-1}, h^{-1} \right)$.
\end{problem}
\begin{solution}
	
	
	Pokażemy, że $\left(h^{-1}\cdot n^{-1}, h^{-1} \right) * \left( n,h \right) = \left( n,h \right) * \left( h^{-1}\cdot n^{-1}, h^{-1} \right) = \left( e_{N}, e_{H} \right) $	

 \textbf{Przypomnienie:} Mamy dowolnie dwie grupy H, N oraz H działa na N przez automorfizmy, wtedy $\forall\left(h,h' \in H\right)\forall\left(n,n' \in N\right)$ zachodzi 
 \begin{enumerate}[label=\alph*)]
 	\item $e_{H}\cdot n = n$,
	\item $h\cdot \left( h' \cdot n \right) = \left( hh' \right)\cdot n$,
	\item $\left( h\cdot n \right)\left( h\cdot n' \right) = h\cdot \left( nn' \right)$
 \end{enumerate}
 Mamy wtedy też $c')\quad h\cdot e_{N} = e_{N}$
	\begin{align*}
		\left( h^{-1}\cdot n^{-1},h^{-1} \right) * \left( n,h \right) =& \left( \bm{\left( h^{-1}\cdot n^{-1} \right)\left( h^{-1}\cdot n \right)},h^{-1}h  \right) \overset{c)}{=} \\
		\left( \bm{h^{-1}\cdot \left( n^{-1}n \right)}, e_{H} \right) =& \left( \bm{h^{-1}\cdot e_{N}}, e_{H} \right) \overset{c')}{=} \left( \bm{e_{N}}, e_{H} \right) \\
	\text{oraz} \\
		\left( n,h \right) * \left( h^{-1}\cdot n^{-1}, h^{-1} \right) =& \left( n\left(\bm{h\cdot \left( h^{-1}\cdot n^{-1}} \right)\right), hh^{-1} \right) \overset{b)}{=} \\
		\left( n\left(\bm{\left( hh^{-1} \right) \cdot n^{-1}} \right), e_{H} \right) =& \left( n\left( \bm{e_{H}\cdot n^{-1}} \right), e_{H}  \right) \overset{a)}{=} \left( n\bm{n^{-1}}, e_{H} \right) = \left( e_{N}, e_{H} \right)  
	.\end{align*}
\end{solution} 
\begin{problem}[2]
	
	Załóżmy, że $\cdot _{1}$ jest działaniem przez automorfizmy grupy $H_1$ na grupie $N_1$. Niech $f: H_1\rightarrow H_2$ oraz $g: N_1\rightarrow N_2$ będą izomorfizami. Definiujemy dla $h_2\in H_2, n_2 \in N_2$
	\[
	h_2 \cdot _{2} n_2 :=g\left( f^{-1}\left( h_1 \right) \cdot_{1}g^{-1}\left( n_2 \right)  \right) 
	.\] 
	\begin{enumerate}[label=\alph*)]
		\item Sprawdzić, że $\cdot _{2}$ jest działaniem przez automorfizmy grupy $H_2$ na grupie $N_2$,
		\item Udowodnić, że funkcja $F: N_1 \rtimes_{\cdot _{1}}H_1 \rightarrow N_2\rtimes_{\cdot _{2}} H_2$ zadana wzorem $F\left( \left( n_1,h_1 \right)  \right) = \left( g\left( n_1 \right) , f\left( h_1 \right)  \right) $ jest izomorfizmem.
	\end{enumerate}
\end{problem} 
\begin{solution}
	
	TODO
\end{solution} 
\begin{problem}[3]
	
	Udowodnić, że $\text{Aut}\left( \Z_2\times \Z_2 \right) \cong \Z_3$
\end{problem} 
\begin{solution}
	
	TODO
\end{solution} 
\begin{problem}[4]
	
	Niech $\varphi: \Z_2\rightarrow \text{Aut}\left( \Z_{n} \right) $ będzie działaniem z zadania 9. z listy 3.. Udowodnić, że $D_{n}\cong \Z_{n}\rtimes \Z_2.$
\end{problem} 
\begin{solution}
	
	TODO
\end{solution} 
\begin{problem}[5]
	
	Udowodnić, że $A_4 \cong \left( \Z_2\times \Z_2 \right) \rtimes \Z_3$.
\end{problem} 
\begin{solution}
	
	TODO
\end{solution} 
\begin{problem}[6]
	
	Niech G będzie podgrupą grupy $S_{\R}$ składającą się z bijekcji \textit{afinicznych}, tzn postaci $x \mapsto ax+b $. Udowodnić, że $G\cong\left( \R,+ \right) \rtimes \left( \R^{*},\cdot  \right) $.
\end{problem} 
\begin{solution}
	
	TODO
\end{solution} 
\begin{problem}[7]
	
	Niech G będzie grupą rzędu 8. Załóżmy, że $g\in G$ jest elementem rzędu 4 i że wszystkie elementy z $G\setminus\langle g \rangle $ są rzędu 4. Dowieść, że $g^2 \in Z\left( G \right) $.
\end{problem} 
\begin{solution}
	
	Ustalmy dowolne $h\in G$. Rozważmy przypadki:
	\begin{enumerate}[label=\arabic*.]
		\item $h\in \langle g \rangle $:

			Wtedy $h = g^{n}$, dla pewnego $0\le n\le 3$. Zatem $hg^2 = g^{n}g^2 = g^{n+2} = g^{2+n} = g^2g^{n} = g^2h$.
		\item $h\not\in \langle g \rangle $ :

			Wtedy $hg^2 = k$ dla pewnego $k \in G$, ponieważ $h\not\in \langle g \rangle $ więc rząd$\left( h \right) = 4$. Zatem $g^2 = h^3k$, a co za tym idzie 
			\begin{align*}
				e = g^{4} =& g^2g^2 = h^3kh^3k\\
				h =& kg^2 \\
				g^2 =& h^3k
			.\end{align*}
			Wtedy 
			\begin{align*}
				g^2h &= \\
				g^2h\bm{e} &= g^2h\bm{h^3kh^3k} = \\
				g^2\bm{hh^3}kh^3k &= g^2\bm{e}kh^3k = \\
				g^2 k\bm{kh^3} &= g^2kg^2 = \\
				\bm{g^2k}g^2 =& \bm{h}g^2\\
					=& hg^2
			.\end{align*}
	\end{enumerate}
\end{solution} 
\begin{problem}[8]
	
	Niech $H_1\unlhd H$ oraz $G_1\unlhd G$. Udowodnić, że $H_1\times G_1 \unlhd H\times G$ oraz (korzystając z zasadniczego twierdzenia o homomorfizmach grup), że 
	\[
		\left( H\times G \right) / \left( H_1\times G_1 \right)\cong \left( H /H_1\right) \times \left( G / G_1 \right)   
	.\] 
\end{problem} 
\begin{solution}
	Ustalmy dowolny $\left( h_1,g_1 \right) \in H_1\times G_1 $ oraz $\left( h,g \right) \in H\times G$.

	Wtedy
	\begin{align*}
		\left( h,g \right) \left( h_1,g_1 \right)\left( h,g \right) ^{-1} =& \left( h,g \right) \left( h_1,g_1 \right)\left( h^{-1},g^{-1} \right) \\ 
		\left( h,g \right) \left( h_1h^{-1}, g_1g^{-1} \right) =& \left( hh_1h^{-1}, gg_1g^{-1} \right)
	.\end{align*}
	Ponieważ $H_1\unlhd H$ oraz $G_1 \unlhd G$, więc $hh_1h^{-1} \in H_1$ oraz $gg_1g^{-1} \in G_1$.

	Zatem $\left( hh_1h^{-1}, gg_1g^{-1} \right) \in \left( H_1, G_1 \right) $, a co za tym idzie
	\[
	H_1\times G_1\unlhd H\times G.
	.\] 

	Ustalmy $\varphi: H\times G \rightarrow \aut\left(\right) $

\end{solution} 
\begin{problem}[9]
	
	Niech $\varphi: G \rightarrow \text{Aut}\left( H \right) $ będzie działaniem. Udowodnić, że następujące warunki są równoważne.
	\begin{enumerate}[label=\alph*)]
		\item Grupa $H\rtimes_{\varphi}G $ jest przemienna.
		\item Grupy $H$ i $G$ są przemienne oraz $\varphi$ jest trywialne.
	\end{enumerate}
\end{problem} 
\begin{solution}
	TODO
\end{solution} 
\begin{problem}[10]
	
	Niech $\Psi: G\rightarrow H$ będzie epimorfizmem. Załóżmy, że istnieje \textit{cięcie} $\Psi$, tzn. homomorfizm $s: H \rightarrow G$, taki że $\Psi \circ s = \id_{H}$. Opisać naturalne działanie $H$ na $\ker(\Psi) $ i udowodnić, że
	\[
	G \cong \ker(\Psi) \rtimes H
	.\] 
\end{problem} 
\begin{solution}
	
	Zauważmy, że $s$ jest \textbf{monomorfizmem}. Istotnie, dla dowolnego $x\in \ker(s), s(x) = e_{G}$. Ponieważ $\left(\Psi\circ s\right)= \id_{H}$, więc
	\[
	e_{H} = \Psi\left( e_{G} \right) = \Psi\left( s\left( x \right)  \right) = \left( \Psi\circ s \right) \left( x \right) = \id_{H}\left( x \right) = x
	.\] 

	Zatem $s[H]$ jest zanurzeniem grupy  $H$ w grupę $G$. Ponadto  $S\left[ H \right] \cap \ker(\Psi) = e_{G}$(bo $s$ jest monomorfizmem). Ponieważ $\ker(\Psi) \unlhd G$, więc $G\cong \ker(\Psi) \rtimes s\left[ H \right] $. Z tego, że $\Psi$ jest epimorfizmem, czyli $\Psi\left[ H \right] = H$, więc $\Psi\rvert_{s\left[ H \right] }$ izomorfizmem $s\left[ H \right] \rightarrow H$. Zatem $G \cong \ker(\Psi) \rtimes s\left[ H \right] \cong \ker(\Psi) \rtimes H$.


		Niech oraz $G' = N'\rtimes H' = \left( \ker(\Psi)\times H, * \right) $. Rozważmy podgrupy $\{e_{G}\}\times H$ oraz $\ker(\Psi) \times \{e_{H}\} $ grupy $G'$. 
	\begin{goall}
		
		
	\end{goall}







	\textbf{Jak wygląda "naturalne" działanie $H$ na $\ker(\Psi) $?}
	
	\begin{goall}[Naturalne działanie H na $\ker(\Psi)$]
	

	\end{goall}


	\begin{figure}[H]
\centering
\begin{tikzpicture}[
    node distance=3cm, % Zmniejszony odstęp poziomy
    % Style
    grupa/.style={draw, ellipse,  align=center, inner sep=0.001pt},
    subgrupa/.style={draw, ellipse, minimum height = 1cm,minimum width = 1.4cm, fill=gray!10, label={below:{\scriptsize$\ker(\Psi)$}}},
    obraz-g/.style={draw, ellipse, minimum height=2.8cm, minimum width = 2.1cm, fill=gray!10, label={above:$H = \Psi[G]$}},  
    obraz-s/.style={draw, ellipse, minimum height=2.8cm, minimum width=2.1cm, fill=blue!10, label={[label distance = 0.25pt]above:{\scriptsize $s[H]$}}},
    punkt/.style={font=\large}
]

% --- Grupa G (Lewa strona) ---
% 1. Komponenty wewnętrzne
\node[subgrupa] (K) {}; % Jądro
\node[obraz-s, above=0.025cm of K] (sH) {}; % Obraz s(H)
\node[punkt] (eG) at (K.north){$\bullet$}; % e_G
% 2. Duży owal G
\node[grupa, fit=(K) (eG) (sH), label={above:$G$}] (G) {};

% 3. Podpis e_G
\node[below=0.025cm of eG] {\scriptsize $e_G$};

% --- Grupa H (Prawa strona) ---
\node[obraz-g, right=of G, minimum width=2.5cm] (H) {};
\node[punkt] (eH) at (H.south) {$\bullet$}; % e_H
\node[below=0.025cm of eH] {\scriptsize $e_H$};

% --- Strzałki (Mapowania) ---
    
% Epimorfizm Psi
\draw[->>, thick] (G.north) to[out=30, in=135] 
    node[midway, above] {$\Psi$} (H.north west);
    
% Cięcie s
\draw[->, thick] (H.west) to[out=180, in=0] 
    node[midway, below] {$s$} (sH.east);
    
% Mapowanie jądra
\draw[->, thick, dashed] (K.north) to[out=0, in=180] (eH);
\draw[->, thick, dashed] (K.south) to[out=0, in=180] (eH);

\end{tikzpicture}
\end{figure}

\end{solution} 
\begin{problem}[11]
	
	Uzasadnić, że w każdym wierszu w klasyfikacji grup rzędu $\le 8$ podanej na wykładzie wypisane są parami nieizomorficzne grupy.
\end{problem} 
\begin{solution}
	
	
\end{solution} 
